\begin{center}
    {\LARGE \textbf{2007无机化学}} \\[2ex]
    {\large 时量180分钟 \quad 满分100分}
\end{center}

\vspace{0.5cm}
\noindent\textbf{注意事项:}
\begin{enumerate}[label=\arabic*., parsep=0pt, itemsep=0pt]
    \item 答题前,考生先将自己的姓名、学号、考场号、座位号填写在答题卡上,将条形码准确粘贴在考生信息条形码粘贴区。
    \item 考试结束后,将本试卷与答题纸一并收回。
    \item 回答各个问题时,将答案写在答题纸上,写在本试卷上无效。
\end{enumerate}

\vspace{0.5cm}
\noindent\textbf{一、完成并配平下列反应的方程式（10 pts.）}

\begin{enumerate}[label=\arabic*., start=1, itemsep=1.5ex]
    \item 向碘水溶液中通入\ce{H2S}。
    \item 向\ce{Na2SO3}溶液中滴加酸性\ce{K2Cr2O7}溶液。
    \item 将酸性\ce{KMnO4}溶液加入\ce{KNO2}溶液。
    \item 向盛有\ce{HgO}的试管中滴加\ce{HI}溶液至过量。
\end{enumerate}

\vspace{0.5cm}
\noindent\textbf{二、解释下列实验现象（40 pts.）}

\begin{enumerate}[label=\arabic*., start=1, itemsep=1.5ex]
    \item 向\ce{CuSO4}溶液中加入\ce{KI}溶液，有黄色沉淀生成；向该黄色沉淀中加入过量\ce{Na2SO3}溶液，沉淀转为白色；若向黄色沉淀中加入过量\ce{Na2S2O3}溶液，沉淀消失，得到无色溶液。
    \item 向\ce{Fe2(SO4)3}溶液中加入\ce{HCl}溶液，溶液变成浅黄色，再加入饱和\ce{K2C2O4}溶液，溶液变为黄绿色。
    \item 白色的\ce{Mn(OH)2}用碘水处理后变为棕褐色，再加入稀硫酸则棕褐色沉淀消失。
    \item 向酸性\ce{K2Cr2O7}溶液中加入\ce{H2O2}，溶液变蓝，过一段时间溶液的蓝色消失。
    \item 绿色\ce{CuCl2 * 2H2O}晶体用煤气灯加热变黑，但蓝色\ce{CuSO4 * 5H2O}晶体用煤气灯加热变白。
\end{enumerate}

\vspace{0.5cm}
\noindent\textbf{三、简要回答下列各题（60 pts.）}

\begin{enumerate}[label=\arabic*., start=1, itemsep=1.5ex]
    \item 给出原子序数为27和45的元素的元素符号、价电子结构。
    \item 分别指出下列化合物所含有的全部化学键类型。
    
    {\centering
    \ce{O3}，\ce{P2O5}，\ce{K3[Fe(CN)6]} \par
    }
    \item 给出下列分子或离子的几何构型和中心原子或离子的杂化类型。
    
    {\centering
    \ce{SO2}，\ce{SOCl2}，\ce{NO2^{+}}，\ce{Ni(CN)4^{2-}}，\ce{Ni(NH3)6^{2+}} \par
    }
    \item 说明原因：\ce{CuCl2 * 2H2O}晶体为绿色而\ce{CuSO4 * 5H2O}晶体为蓝色。
    \item 用分子轨道理论比较物质的稳定性：
    
    {\centering
    $\ce{O2}$与$\ce{O2^{+}}$，$\ce{NO}$与$\ce{NO^{+}}$ \par
    }
    \item 画出下列配合物所有可能存在的异构体：
    
    {\centering
    \ce{Ni(en)2(NH3)2}，\ce{PtCl2(OH)(NH3)3} \par
    }
    \item 比较各对化合物热稳定性并说明原因：
    
    {\centering
    $\ce{HNO3}$与$\ce{HNO2}$，$\ce{Li2CO3}$与$\ce{K2CO3}$ \par
    }
    \item 已知：\\
    \ce{Fe^{3+} + e^{-} = Fe^{2+}} \quad $E^\ominus = 0.77\ \text{V}$ \\
    \ce{Fe(OH)3 + e^{-} = Fe(OH)2 + OH^{-}} \quad $E^\ominus = -0.56\ \text{V}$ \\
    \ce{Fe(CN)6^{3-} + e^{-} = Fe(CN)6^{4-}} \quad $E^\ominus = +0.36\ \text{V}$
    \\
    比较：\ce{Fe(OH)3}和\ce{Fe(OH)2}溶度积常数大小；\ce{Fe(CN)6^{3-}}和\ce{Fe(CN)6^{4-}}稳定常数大小
\end{enumerate}

\vspace{0.5cm}
\noindent\textbf{四、制备与分离（15 pts.）}

\begin{enumerate}[label=\arabic*., start=1, itemsep=1.5ex]
    \item 用反应方程式（不用配平）表示下列转化过程：
    
    \ce{Cu -> CuCl2 * 2H2O -> K2[Cu(C2O4)2] * 2H2O}
    \item 用反应方程式（不用配平）表示：\ce{MnO2}为主要原料制备水合晶体\ce{MnSO4 * 7H2O}
    \item 设计方案分离溶液中的：\ce{Cr^{3+}}、\ce{Al^{3+}}、\ce{Fe^{3+}}、\ce{Ag^{+}}、\ce{Ni^{2+}}
\end{enumerate}

\vspace{0.5cm}
\noindent\textbf{五、鉴别与判断（25 pts.）}

\begin{enumerate}[label=\arabic*., start=1, itemsep=1.5ex]
    \item 用二种方法鉴别下列各对物质：
    
    {\centering
    \ce{SbCl3}和\ce{BiCl3}，\ce{CuO}和\ce{MnO2} \par
    }
    \item 绿色水合硫酸盐晶体\ce{A}溶于水后加入适量氨水生成绿色沉淀\ce{B}。氨水过量时，\ce{B}溶解生成蓝色溶液\ce{C}，\ce{B}不溶于\ce{NaOH}溶液，但在强碱性条件下与氯水作用生成棕黑色沉淀\ce{D}。将\ce{D}与\ce{MnSO4}混合后加入稀硫酸充分反应，再加入\ce{NaOH}溶液得到绿色沉淀\ce{B}和紫色溶液\ce{E}。请给出\ce{A} $\sim$ \ce{E}的化学式。
    \item 化合物\ce{A}为绿色晶体。\ce{A}溶于水后加入\ce{AgNO3}溶液生成白色沉淀\ce{B}，\ce{B}不溶于\ce{HNO3}溶液但溶于\ce{NaOH}溶液，但溶于氨水。向\ce{A}的溶液中加入氨水生成灰蓝色沉淀\ce{C}。\ce{C}溶于\ce{NaOH}溶液得绿色溶液\ce{D}。向溶液\ce{D}中加入\ce{H2O2}则溶液变黄，说明有溶液\ce{E}生成。向溶液\ce{E}中加入\ce{BaCl2}溶液生成黄色沉淀\ce{F}。\ce{F}溶于\ce{HNO3}溶液而不溶于\ce{NaOH}溶液。用稀\ce{H2SO4}处理沉淀\ce{F}得到白色沉淀和橙色溶液。请给出\ce{A} $\sim$ \ce{F}的化学式。
\end{enumerate}